\section*{Chapitre \ref{ch:ml:low-latency-inference} --- Inférence à faible latence}

\begin{solution}{exo:ml:low-latency-inference:1}
L'arrondi au multiple le plus proche de $s$ commet une erreur d'au plus $s/2 = \max|w|/254$~: 0.39~\% de la plus grande
pondération, et une part bien plus grande d'une petite pondération.
\end{solution}

\begin{solution}{exo:ml:low-latency-inference:2}
$5\,738/544\,349 = 0.011$~: le 99\up{e}~centile de la forêt compilée vaut environ un quatre-vingt-quinzième de celui de l'appel à la bibliothèque.
\end{solution}

\begin{solution}{exo:ml:low-latency-inference:3}
La prédiction d'une forêt est une somme de valeurs de feuilles choisies par comparaisons~: la version compilée effectue les mêmes comparaisons et additionne les mêmes doubles dans le même ordre.
Le calcul en virgule flottante d'un réseau dépend de l'ordre des sommes dans chaque produit, des multiplications-additions fusionnées (FMA) et de la
vectorisation, qui diffèrent d'une bibliothèque ou d'un compilateur à l'autre~; la parité n'est alors qu'une tolérance, sauf si
l'arithmétique est entière, comme dans le réseau int8.
\end{solution}

\begin{solution}{exo:ml:low-latency-inference:4}
$0.0123 = 0.7872\times2^{-6}$~; $M = \mathrm{round}(0.7872\times2^{31}) = 1\,690\,499\,128$ et un décalage total de $6 + 31 = 37$~: le
produit est $\mathrm{acc}\times M$ décalé de 37 bits vers la droite. Les entiers rendent le résultat identique sur toutes les plateformes et évitent
la virgule flottante sur le matériel qui n'en dispose pas, ou où elle est plus lente.
\end{solution}

\begin{solution}{exo:ml:low-latency-inference:5}
À huit bits, l'arrondi est trop fin pour nuire, si bien que l'\omterm{def:ml:representation-learning:ssl}{ajustement fin} (fine-tuning) ne fait que éloigner le modèle de son optimum entraîné~; à quatre bits,
l'arrondi est assez grossier pour dégrader le modèle post-entraînement, et un modèle entraîné à l'anticiper récupère plus de la moitié de la
perte (0.164 contre 0.143, en virgule flottante 0.182).
\end{solution}

\begin{solution}{exo:ml:low-latency-inference:6}
Diviser le temps d'un lot par sa taille donne l'inverse du débit, et non la latence d'une requête~: cela masque l'attente de la première
requête, la queue de distribution et tout ce qui survient une fois par appel. Mesurer les appels isolés, au fil de leur arrivée, et rapporter la distribution.
\end{solution}

\begin{solution}{exo:ml:low-latency-inference:7}
Des échelles par sortie donnent à chaque ligne de pondérations sa propre grille, si bien que les petites lignes ne sont pas aplaties par une grande pondération ailleurs~: à quatre
bits, le coefficient d'information (IC) passe de 0.143 à 0.153 (en virgule flottante 0.182), soit environ un tiers de la perte récupérée
sans aucun réentraînement. À huit bits, il vaut 0.182, comme précédemment.
\end{solution}

\begin{solution}{exo:ml:low-latency-inference:8}
Avec $x_i = s_xX_i$ et $w_i = s_wW_i$, $\sum_iw_ix_i + b = s_xs_w\big(\sum_iW_iX_i + b/(s_xs_w)\big)$~: stocker $B =
\mathrm{round}(b/(s_xs_w))$ à l'échelle de l'accumulateur permet d'ajouter le biais à la somme entière. L'entier de sortie est
$Y = \mathrm{round}\big((s_xs_w/s_y)(\sum_iW_iX_i + B)\big)$~: le multiplicateur est $s_xs_w/s_y$, appliqué comme un produit en virgule fixe
et un décalage.
\end{solution}

\begin{solution}{pb:ml:low-latency-inference:1}
\textbf{Partie I.}
\begin{enumerate}
\item Une forêt LightGBM de 300 arbres (IC 0.224) et un réseau $16\to32\to16\to1$ (IC 0.182).
\item Sous forme de tableaux plats de nœuds (parcourus par la même boucle en Python, C++ et Rust) et sous forme de branches C++ générées.
\item Arbres~: les mêmes comparaisons et le même ordre d'additions en double~; réseau~: arithmétique entière.
\item Les 4\,200 embranchements et les 4\,500 feuilles de la forêt tiennent dans des tableaux de quelques centaines de kilo-octets de texte, et bien moins en binaire~; le
réseau pèse 1\,236 octets en int8, contre 4\,356 en float32.
\end{enumerate}
\textbf{Partie II.}
\begin{enumerate}[resume]
\item Médiane et 99\up{e}~centile (µs)~: LightGBM 265 et 544~; forêt Python 300 et 397~; réseau PyTorch 21 et 44~; réseau int8 NumPy
24 et 66~; boucle de forêt C++ 6.2 et 13.2~; branches C++ 2.6 et 5.7~; réseau int8 C++ 0.47 et 1.1~; forêt Rust 7.8
et 17.8~; réseau int8 Rust 0.58 et 1.9.
\item L'interpréteur, les vérifications d'arguments, les conversions de la ligne vers et depuis les tableaux et la mécanique du pool de threads~; les arbres
eux-mêmes ne prennent que quelques microsecondes.
\item Les seuils et les cibles de branchement sont dans le code, et non chargés depuis la mémoire, et le compilateur agence les chemins d'exécution.
\item Environ un quart, pour les contrôles de bornes à chaque accès aux tableaux~; ce coût est supprimable en validant le modèle au chargement.
\end{enumerate}
\textbf{Partie III.}
\begin{enumerate}[resume]
\item Rien de mesurable après l'entraînement (0.182)~; \omterm{def:ml:representation-learning:ssl}{ajustement fin} conscient de la \omterm{def:ml:low-latency-inference:quant}{quantification}~: 0.179.
\item La \omterm{def:ml:low-latency-inference:quant}{quantification post-entraînement} tombe à 0.143, l'\omterm{def:ml:low-latency-inference:quant}{entraînement conscient de la quantification} tient à 0.164, les échelles par canal à 0.153.
\item De 273~µs par appel pour une seule ligne à 3.1~ms pour 1\,024 lignes~: 3.0~µs par ligne.
\item Pour noter de nombreuses requêtes sans contrainte de latence (un univers chaque seconde)~; non pour une décision sur un instrument isolé.
\end{enumerate}
\textbf{Partie IV.}
\begin{enumerate}[resume]
\item \emph{Deux microsecondes pour une forêt.} Le 99\up{e}~centile de la forêt C++ générée vaut 5.7~µs, contre 544~µs pour
l'appel à la bibliothèque, soit un rapport d'environ 1 pour 95 (2.6 contre 265~µs à la médiane)~; la \omterm{def:ml:low-latency-inference:quant}{quantification} int8 ne coûte rien de mesurable au réseau après
l'entraînement (0.182) et 0.003 avec l'\omterm{def:ml:representation-learning:ssl}{ajustement fin} conscient de la \omterm{def:ml:low-latency-inference:quant}{quantification} (0.179)~; à quatre bits, 0.040 et 0.018.
\item La forêt en C++ généré (2.6~µs en médiane, 5.7 à la queue) si son signal supplémentaire en vaut la peine, ou le réseau int8
à une demi-microseconde.
\item Le modèle source, l'artefact compilé et son empreinte (hash), le compilateur et ses options, et le rapport de parité.
\item Cœurs isolés, aucun autre travail, caches chauds comme en production, requêtes isolées au rythme de production, et queues de distribution sur de longues
exécutions, en mesurant le surcoût propre à la mesure elle-même.
\item Dix fois plus de code~: le source généré ne tiendra peut-être plus dans le cache d'instructions, et des tableaux plats ou une évaluation
vectorisée peuvent l'emporter.
\item Lorsque le budget est inférieur à la microseconde et fixe, et que le modèle est assez petit et stable pour justifier un cycle matériel.
\item Sous la forme d'un écart entre les variables que reçoit le modèle compilé et celles qu'a vues le modèle de recherche, invisible aux
tests de parité du noyau~: la parité doit aussi être vérifiée sur des entrées réelles.
\item Surtout autour du modèle, et non dans celui-ci.
\end{enumerate}
\end{solution}

\begin{solution}{iq:ml:low-latency-inference:1}
Exporter les arbres et les compiler (branches générées ou tableaux plats en C++), garder le processus à chaud, éviter les allocations et
les conversions sur le hot path, vérifier la parité avec la bibliothèque et mesurer les queues une requête à la fois.
\iqlookfor{compilation, un hot path propre, parité et mesure des queues.}
\end{solution}

\begin{solution}{iq:ml:low-latency-inference:2}
Les valeurs réelles sont des entiers multipliés par une échelle (plus un point zéro dans les schémas asymétriques)~; les produits
s'accumulent en int32~; l'accumulateur de chaque couche est ramené aux entiers de la couche suivante par un multiplicateur en virgule fixe~; les biais sont stockés à l'échelle de l'accumulateur.
\iqlookfor{échelles, accumulateurs et requantification.}
\end{solution}

\begin{solution}{iq:ml:low-latency-inference:3}
Aucune allocation sur le tas, aucun verrou, aucun appel système, aucune exception ni appel virtuel~; boucles bornées~; données
disposées pour le cache~; tout est chargé et validé au démarrage.
\iqlookfor{les règles usuelles du hot path, avec leurs raisons.}
\end{solution}

\begin{solution}{iq:ml:low-latency-inference:4}
Vecteurs de test partagés avec des sorties de référence issues du framework d'entraînement~; comparaison exacte pour les arbres et les réseaux entiers, tolérances énoncées
pour les réseaux flottants~; empreinte (hash) de l'artefact compilé enregistrée avec le modèle source~; contrôles aussi sur les entrées réelles.
\iqlookfor{vecteurs partagés, exactitude quand c'est possible, enregistrement.}
\end{solution}

\begin{solution}{iq:ml:low-latency-inference:5}
Elle aide lorsque le surcoût par appel domine et que les requêtes peuvent attendre~; elle nuit lorsque chaque requête a une échéance, car la première
requête d'un lot attend la dernière.
\iqlookfor{amortissement du surcoût face au temps d'attente.}
\end{solution}

\begin{solution}{iq:ml:low-latency-inference:6}
Quantifier après l'entraînement et mesurer la métrique qui compte~; si la perte est négligeable, s'arrêter~; sinon, essayer une granularité plus fine
(par canal) et, si besoin, un \omterm{def:ml:representation-learning:ssl}{ajustement fin} conscient de la \omterm{def:ml:low-latency-inference:quant}{quantification}, puis comparer sur des données hors échantillon.
\iqlookfor{mesurer d'abord, puis passer à plus fin.}
\end{solution}
